\chapter{Conclusioni}
\label{cap:concl}

Questa tesi ha affrontato il problema del \emph{source seeking} marino --- la
localizzazione autonoma della sorgente di un inquinante a partire dalle letture locali
di concentrazione --- confrontando due famiglie di approcci nell'ambito del progetto
HYDRAS, su un dataset di simulazioni idrodinamiche realistiche della baia di Cecina.

\section{Risultati principali}

Ripercorriamo il cammino della tesi dall'inizio alla fine, raccogliendone i risultati
salienti.

\paragraph{Il problema e i dati.}
Il punto di partenza \`e un problema concreto: localizzare in modo autonomo la sorgente di
uno sversamento inquinante dalle sole letture locali di concentrazione, in un ambiente
costiero dove la dispersione \`e governata da una dinamica complessa di correnti, vento e
turbolenza. Nell'ambito del progetto HYDRAS, il compito \`e stato affrontato su un dataset
di simulazioni idrodinamiche ad alta risoluzione (MIKE21) della baia di Cecina
(Capitolo~\ref{cap:dati}), forzate dai dati meteo-oceanografici di CMEMS ed ERA5: 132
sorgenti equispaziate, ciascuna simulata sotto quattro scenari di vento (V0--V3) che
generano plume dalle morfologie marcatamente diverse. Per una valutazione rigorosa, 26
sorgenti (SRC107--132) sono state riservate alla sola inferenza --- mai viste in
addestramento --- e ogni simulazione \`e stata suddivisa in tre finestre temporali,
definendo il protocollo \emph{held-out} su cui entrambi gli approcci sono stati misurati a
parit\`a di condizioni.

\paragraph{L'ambiente condiviso.}
Su questi dati \`e stato costruito un unico ambiente di simulazione
(Capitolo~\ref{cap:ambiente}), condiviso dai due approcci: una squadra \emph{master--slave}
in cui un agente master decide il movimento e otto slave, disposti in corona lungo le
direzioni cardinali e diagonali, campionano la concentrazione e gliela riportano. L'ambiente
definisce lo spazio d'azione discreto, l'\emph{action masking} che impedisce le collisioni
con la costa, la procedura di inizializzazione sul plume e un segnale di reward. Far operare
i due metodi sullo stesso ambiente --- stessa dinamica, stessi vincoli, stesse osservazioni
--- \`e ci\`o che rende il confronto pulito: l'unica variabile \`e il modulo che, ricevuta
l'osservazione, sceglie l'azione.

\paragraph{Il metodo del gradiente.}
Il primo modulo \`e il \textbf{Field Climbing Method} (Capitolo~\ref{cap:fcm}), la baseline
gradient-based. Nella versione a passo fisso raggiunge un success rate globale di appena il
\textbf{32{,}9\%}, limitato dalla sua natura puramente reattiva. L'introduzione
dell'ottimizzatore \textbf{Adam} per il calcolo del passo --- con momentum sul gradiente
stimato e normalizzazione per la varianza storica --- lo porta fino al \textbf{63{,}5\%}
($lr = 40\,\text{m}$), con un guadagno di oltre trenta punti: il momentum smorza le
oscillazioni dovute al rumore del sensore. Il passo ottimale non \`e per\`o il pi\`u ampio
--- oltre i $40\,\text{m}$ falcate troppo lunghe scavalcano la struttura del plume --- e i
valori pi\`u performanti implicano velocit\`a superiori a quella nominale del veicolo. Al di
l\`a di questi affinamenti, il limite resta \emph{strutturale}: negli scenari a plume diffuso
il gradiente locale degenera e il metodo, privo di memoria e di contesto, si blocca.

\paragraph{L'apprendimento per rinforzo.}
Il secondo modulo \`e un agente addestrato con \textbf{PPO} (Capitolo~\ref{cap:rl}). Uno
studio di ablation sulla funzione di reward ha mostrato che la variante pi\`u semplice
(\texttt{minimal}), basata sul solo gradiente chimico come segnale direzionale, ottiene il
success rate pi\`u alto (\textbf{96{,}9\%}): i reward di allineamento a vento e corrente si
sono rivelati addirittura controproducenti, perch\'e fondati sull'ipotesi --- non valida in
ambiente costiero --- di una sorgente sempre controvento. Rendere poi la \textbf{velocit\`a
una scelta dell'agente} (spazio $\text{Discrete}(8\times5)$, addestrato a catena alzando
progressivamente $v_{\max}$) ha portato il success rate al \textbf{$99{,}7$--$99{,}9\%$} e
ridotto di oltre quattro volte il tempo al successo (da $\sim$20 a $\sim$5 minuti simulati).
Su questa catena, la \textbf{doppia corona} di sensori si \`e rivelata un vantaggio netto a
partire da $v_{\max} = 2\,\text{m/s}$ --- stessa affidabilit\`a e $25$--$34\%$ di passi in
meno --- correggendo la lettura di un primo confronto in cui, addestrata da zero, sembrava
dannosa. Approfondimenti mirati hanno infine mostrato che l'agente sa \emph{manovrare}
finemente attorno alla sorgente (test a raggio di successo ridotto) e che l'esito di un
episodio \textbf{non dipende dalla posizione di partenza} (analisi delle \emph{spawn map}),
ma solo dalle condizioni idrodinamiche dello scenario.

\paragraph{Il confronto.}
Il confronto diretto \`e netto: il modello RL supera il metodo del gradiente su ogni scenario.
Poich\'e entrambi dispongono di una versione a velocit\`a variabile, il confronto pi\`u equo
\`e quello \textbf{a parit\`a di velocit\`a massima}, che neutralizza il vantaggio cinematico
di un passo pi\`u lungo: a ogni livello di velocit\`a il PPO domina l'FCM di
\textbf{$36$--$42$ punti} di success rate, con il divario massimo proprio negli scenari a
plume diffuso, dove il gradiente locale degenera. Esteso all'intero range operativo (da
$0{,}1$ a $5\,\text{m/s}$), il quadro non cambia: il PPO supera il $98\%$ gi\`a a
$1{,}2\,\text{m/s}$ e si stabilizza attorno al $99\%$, mentre l'FCM resta confinato attorno al
$55$--$63\%$ alle velocit\`a operative; l'unico regime in cui il divario si assottiglia \`e
quello delle velocit\`a pi\`u basse, dove un deficit puramente fisico di propulsione --- il
veicolo pi\`u lento delle correnti da risalire --- penalizza ogni strategia.

\paragraph{L'interpretazione della politica.}
Stabilito \emph{che} la politica funziona, il Capitolo~\ref{cap:xai} ne ha indagato il
\emph{perch\'e}. Tramite ablation causale di gruppo, valori di Shapley esatti e sonde
comportamentali \`e emerso un ritratto coerente: la politica \textbf{non} \`e un risalitore di
gradiente locale (muove verso il sensore pi\`u intenso solo il $6\%$ delle volte in pi\`u di
quanto predica il suo stesso bias direzionale), ma si regge sull'\textbf{integrazione
temporale} del campo chimico --- la storia della corona di sensori --- unita alla memoria di
dove si \`e andati. \`E emersa inoltre una divergenza istruttiva fra \emph{ci\`o che la rete
guarda} e \emph{ci\`o che le serve}: la corrente marina \`e la grandezza pi\`u ``guardata''
eppure la sua rimozione costa pochissimo, mentre la storia della corona \`e il solo blocco
davvero insostituibile. Un esperimento di riaddestramento \emph{senza} formazione ne ha
infine quantificato il valore: la formazione vale circa \textbf{$6{,}4$ punti} di success
rate, ma \`e in parte sostituibile da una strategia temporale appresa --- l'agente singolo
riaddestrato resta al $93{,}3\%$, circa trentacinque punti sopra l'FCM che pure dispone
dell'intera formazione. Ci\`o che separa il successo dal fallimento, in definitiva, non \`e il
sensore ma la \textbf{strategia appresa}.

\section{Sviluppi futuri}

La direzione di sviluppo pi\`u naturale e ambiziosa \`e il passaggio da un sistema a
\emph{singolo decisore} a un \textbf{sistema multi-agente cooperativo ``vero''}.
Nell'architettura master--slave adottata in questa tesi un solo agente (il master) decide
il movimento, mentre gli otto slave sono sensori passivi in formazione rigida: di fatto il
problema resta a singolo agente. La visione complessiva del progetto HYDRAS prevede invece
una squadra di $n$ robot autonomi, ciascuno dotato di una \emph{propria} politica di
navigazione, che collaborano --- coordinando l'esplorazione, ripartendosi il dominio e
scambiandosi informazioni --- per localizzare la sorgente pi\`u rapidamente e con maggiore
robustezza di un sistema centralizzato.

\paragraph{L'algoritmo MAPPO.} Per addestrare una tale squadra, l'estensione naturale del
PPO impiegato in questa tesi \`e \textbf{MAPPO} (\emph{Multi-Agent PPO})~\citep{mappo}, che
adatta il PPO al contesto multi-agente cooperativo secondo il paradigma \textbf{CTDE}
(\emph{Centralized Training, Decentralized Execution}). L'idea di fondo \`e separare ci\`o
che serve in addestramento da ci\`o che serve in esecuzione:
\begin{itemize}[leftmargin=1.4em, itemsep=3pt]
    \item ogni agente possiede una \textbf{politica decentralizzata} (attore)
          $\pi_\theta(a_i \mid o_i)$ che sceglie l'azione a partire dalla \emph{sola}
          osservazione locale $o_i$ (le proprie letture di concentrazione, la propria
          cinematica). In esecuzione, quindi, ogni robot \`e autonomo: naviga senza un
          coordinatore centrale e senza comunicazione a banda piena;
    \item in addestramento, un \textbf{critico centralizzato} (\emph{value function})
          $V_\phi(s)$ riceve informazione \emph{globale} sullo stato congiunto della
          squadra --- non disponibile ai singoli attori --- e ne stima il valore atteso.
          \`E questo critico ``onnisciente'' a stabilizzare l'apprendimento: valuta le
          azioni di ciascun agente tenendo conto di ci\`o che fanno gli altri, mitigando la
          \emph{non-stazionariet\`a} che affligge gli approcci in cui ogni agente apprende
          in modo del tutto indipendente (IPPO, \emph{Independent PPO}).
\end{itemize}
Il critico viene usato \emph{solo} durante il training e scartato in esecuzione, dove
sopravvivono le sole politiche locali --- da cui il nome CTDE. La scelta di cosa fornire in
ingresso al critico \`e un punto delicato: gli autori mostrano che uno stato
\emph{agent-specific} (lo stato globale dell'ambiente concatenato all'osservazione locale
dell'agente, con le feature ridondanti rimosse) rende meglio sia della semplice
concatenazione di tutte le osservazioni sia del solo stato globale dell'ambiente.

Sul piano meccanico MAPPO conserva l'impianto del PPO di questa tesi --- obiettivo
\emph{clipped}, stima del vantaggio con GAE, normalizzazione del vantaggio --- applicato
per-agente, con due accorgimenti che si rivelano decisivi. Il primo \`e la
\textbf{condivisione dei parametri} tra agenti omogenei: attori e critici usano gli stessi
pesi per tutti i robot, il che moltiplica i dati utili e migliora l'efficienza campionaria
quando i veicoli sono identici. Il secondo \`e la \textbf{normalizzazione del value target}
tramite media e deviazione standard correnti, che evita l'instabilit\`a quando i ritorni
variano molto tra episodi. Gli autori raccomandano inoltre poche epoche PPO per
aggiornamento (5--15, meno nei task difficili) e al pi\`u due mini-batch, per contenere la
deriva delle politiche fra un update e l'altro. Il risultato principale del loro lavoro \`e
che, \textbf{all'aumentare del numero di agenti}, il vantaggio del critico centralizzato di
MAPPO cresce nettamente rispetto agli approcci indipendenti: \`e proprio nei regimi a molti
agenti --- quelli di interesse per una squadra numerosa di robot marini --- che l'algoritmo
d\`a il meglio.

\paragraph{Adattamento al source seeking.} Calato nel nostro dominio, ciascun agente
sarebbe un veicolo autonomo con la propria corona di sensori; il critico centralizzato, in
addestramento, potrebbe accedere alle posizioni e alle letture dell'\emph{intera} squadra
(e, in simulazione, all'intero campo MIKE21), mentre in mare ogni veicolo naviga sulla sola
informazione locale e su una comunicazione sparsa con i vicini. Rispetto al master--slave
qui adottato, ci\`o trasformerebbe gli slave da sensori passivi in decisori attivi, aprendo
a strategie di copertura e di ripartizione del dominio precluse a un singolo agente ---
esattamente l'architettura \emph{multi-agente distribuita} che costituisce l'obiettivo di
lungo periodo del progetto HYDRAS.

\medskip
Restano infine aperte altre direzioni complementari:
\begin{itemize}
    \item \textbf{Validazione su veicolo reale}: trasferire la politica appresa in
          simulazione a un AUV fisico, affrontando il \emph{sim-to-real gap} dovuto al
          rumore reale dei sensori e alla dinamica del veicolo;
    \item \textbf{Generalizzazione geografica}: addestrare e validare gli agenti su
          domini marini diversi dalla baia di Cecina, per valutare la trasferibilit\`a
          delle strategie apprese;
    \item \textbf{Architetture ricorrenti}: sostituire la policy MLP con reti dotate di
          memoria esplicita (LSTM, Transformer) per modellare in modo pi\`u ricco la
          dinamica temporale del plume.
\end{itemize}

I risultati ottenuti confermano il potenziale del Reinforcement Learning come strumento
per la navigazione autonoma in ambienti marini complessi, e quantificano il contributo
dell'apprendimento rispetto agli approcci gradient-based classici.
